\chapter{Exercicis proposats}
\label{cap:exercicis-proposats}

En aquest capítol es proposen exercicis per treballar de manera autònoma els continguts de la unitat. Cada exercici inclou una indicació breu i la solució corresponent.

\section{Definició i operacions amb polinomis}

\begin{exercici}
Escriu un polinomi en $x$, de quart grau, complet i ordenat en
potències decreixents, amb (a) coeficients enters (b) coeficients
racionals i (c) coeficients irracionals
\end{exercici}
\begin{indicacio}
Recorda que un polinomi complet de quart grau ha de contenir termes de grau 4, 3, 2, 1 i 0. Tria els coeficients dins del conjunt indicat.
\end{indicacio}

\begin{solucio}
(a) $x^{4}-5x^{3}+2x^{2}-x+1$, (b) $2x^{4}-\frac{3}{5}x^{3}+
\frac{1}{4}x^{2}-6x+\frac{2}{3}$, (c) $\sqrt{2}x^{4}-3\sqrt{2}x^{3}+\sqrt{3}
x^{2}+3\sqrt{5}x+\sqrt{7}$
\end{solucio}

\begin{exercici}
Per a quin valor de $a$ el polinomi $
p(x)=(a-2)x^{2}+2(a^{2}-4)x+8a-16$ es redueix el polinomi zero?
\end{exercici}
\begin{indicacio}
Un polinomi és el polinomi zero si tots els seus coeficients són zero. Iguala’ls a zero dona les condicions sobre el paràmetre.
\end{indicacio}

\begin{solucio}
$a=2$
\end{solucio}

\begin{exercici}
Quins valors han de tenir $a$ i $b$ perquè els
polinomis $p(x)=x^{3}+(b+2)x^{2}-7$ i $q(x)=x^{3}+ax^{2}+(a+b)x-7$ siguin
iguals?
\end{exercici}
\begin{indicacio}
Iguala els coeficients dels termes del mateix grau o aplica el teorema del residu, segons correspongui.
\end{indicacio}

\begin{solucio}
$a=1$ i $b=-1$
\end{solucio}

\begin{exercici}
Donats els polinomis següents:
\[
\begin{array}{l}
a(x)=-2x^{2}+6x-5 \\
b(x)=3x^{2}+1 \\
c(x)=x^{3}-6x+9 \\
d(x)=\frac{1}{2}x^{3}-3x^{2}-x-\frac{1}{4} \\
e(x)=\frac{1}{3}x^{3}+\frac{1}{2}x^{2}+x
\end{array}
\]
(a) calcula: (i) $a(x)+b(x)+c(x)$, (ii) $3\cdot a(x)-2\cdot b(x)+5\cdot e(x)$, (iii) $a(x)\cdot \left[ b(x)+c(x)\right] $, (iv) $\left[ 3\cdot
d(x)-2\cdot e(x)\right] \left[ a(x)-b(x)\right] $, i (b) troba el polinomi $
p(x)$ que compleix la igualtat proposada: (i) $a(x)+p(x)=e(x)$, (ii) $4\cdot
p(x)=a(x)$, (iii) $5\cdot p(x)-2\cdot a(x)=b(x)+c(x)$.
\end{exercici}
\begin{indicacio}
Fes les operacions indicades, ordena els termes per graus i redueix termes semblants al final.
\end{indicacio}

\begin{solucio}
(a) (i) $x^{3}+x^{2}+5$, (ii) $\frac{5}{3}x^{3}-\frac{19}{2}
x^{2}+23x-17$, (iii) $-2x^{5}+25x^{3}-71x^{2}+90x-50$, (iv) $-\frac{25}{6}
x^{5}+55x^{4}-40x^{3}+\frac{135}{4}x^{2}+\frac{51}{2}x+\frac{9}{2}$, (b) (i)
$\frac{1}{3}x^{3}-\frac{3}{2}x^{2}-5x+5$, (ii) $\frac{1}{2}x^{2}+\frac{3}{2}
x-\frac{5}{4}$ i (iii) $\frac{1}{5}x^{3}-\frac{1}{5}x^{2}+\frac{6}{5}x$
\end{solucio}

\begin{exercici}
Troba el valor numèric de $p(x)=8x^{4}-5x^{2}+5$ per a $x=\sqrt{2}$
i per a $x=-\sqrt{2}$. Què observes?
\end{exercici}
\begin{indicacio}
Substitueix el valor de la variable i observa si només apareixen potències parelles o imparelles.
\end{indicacio}

\begin{solucio}
$p(\sqrt{2})=p(-\sqrt{2})=27$. El valor numèric és el mateix perquè totes les potències de $x$ són parelles.
\end{solucio}

\begin{exercici}
Troba el valor numèric de $p(x)=3x^{5}-5x^{3}+8x$ per a $x=0$, $
x=1 $i per a $x=-1$. Què observes?
\end{exercici}
\begin{indicacio}
Iguala els coeficients dels termes del mateix grau o aplica el teorema del residu, segons correspongui.
\end{indicacio}

\begin{solucio}
$p(0)=0$, $p(1)=6$ i $p(-1)=-6$. Els dos donarrers valors són oposats perquè totes les potències de $x$ són imparelles.
\end{solucio}

\begin{exercici}
Fes els càlculs indicats i redueix termes
semblants:

\begin{enumerate}
\item[a)] $(-2x^{3}-5x+7)\cdot (x^{2}-3x)-(-2x^{2}+x+1)^{2}$

\item[b)] $2\cdot \left[ 3-2x+4\cdot (-3x^{2}-5x+7)\right] -\left[ 6x-2\cdot
(3-2x)\right] ^{2}$

\item[c)] $(x-2)^{3}-(x+3)(x-1)(x+5)-(x-1)(x^{2}+1)-(x-1)^{3}$

\item[d)] $(x-3)(x+3)(x^{2}+9)$

\item[e)] $(3x-4)(2x^{2}+x-6)-(3x-4)(x^{2}+7x+11)-x^{2}(3x-4)$
\end{enumerate}
\end{exercici}
\begin{indicacio}
Fes les operacions indicades, ordena els termes per graus i redueix termes semblants al final.
\end{indicacio}

\begin{solucio}
\[
\begin{array}{ll}
a) & -2x^{5}+2x^{4}-x^{3}+25x^{2}-23x-1 \\
b) & -124x^{2}+76x+26 \\
c) & -2x^{3}-9x^{2}+x+9 \\
d) & x^{4}-81 \\
e) & -18x^{2}-27x+68
\end{array}
\]
\end{solucio}

\begin{exercici}
Transforma les expressions següents en productes:

\begin{enumerate}
\item[a)] $x^{2}-x+\frac{1}{4}$

\item[b)] $(x+a)(x-a)+(x+a)^{2}$

\item[c)] $x^{2}-2ax+a^{2}-b^{2}$

\item[d)] $(x+a-b)^{2}-(x-a+b)^{2}$

\item[e)] $x^{2}-4x+4-(x-2)(x+7)$

\item[f)] $\frac{1}{3}(x-3)(2x-5)^{2}-\frac{1}{3}(x-3)(x+2)^{2}$
\end{enumerate}
\end{exercici}
\begin{indicacio}
Busca productes notables, diferències de quadrats i factors comuns abans de desenvolupar completament.
\end{indicacio}

\begin{solucio}
\[
\begin{array}{ll}
a) & (x-1/2)^{2} \\
b) & 2x(x+a) \\
c) & (x-a+b)(x-a-b) \\
d) & 4x(a-b) \\
e) & -9(x-2) \\
f) & (x-3)(x-1)(x-7)
\end{array}
\]
\end{solucio}

\begin{exercici}
Determina $m$ i $n$ perquè els polinomis $p(x)=x^{3}+2x^{2}+mx+6$ i
$q(x)=(x+n)\cdot (x^{2}+nx+6)$ siguin iguals.
\end{exercici}
\begin{indicacio}
Iguala els coeficients dels termes del mateix grau o aplica el teorema del residu, segons correspongui.
\end{indicacio}

\begin{solucio}
$m=7$ i $n=1$
\end{solucio}

\section{Potència d'un binomi i fórmula de Newton}

\begin{exercici}
Desenvolupa les potències indicades, usant la fórmula del
binomi de Newton:

\begin{enumerate}
\item[a)] $(2x-3)^{4}$

\item[b)] $\left( x-\frac{1}{x}\right) ^{3}$

\item[c)] $(a^{2}b+3ab^{2})^{5}$

\item[d)] $\left( x^{2}-\frac{2}{x}\right) ^{4}$
\end{enumerate}
\end{exercici}
\begin{indicacio}
Usal terme general del binomi de Newton: el terme corresponent depèn del coeficient binomial i de les potències dels dos termes del binomi.
\end{indicacio}

\begin{solucio}
a) $16x^{4}-96x^{3}+216x^{2}-216x+81$; b) $x^{3}-3x+\frac{3}{x}
-\frac{1}{x^{3}}$; c) $
a^{10}b^{5}+15a^{9}b^{6}+90a^{8}b^{7}+270a^{7}b^{8}+405a^{6}b^{9}+243a^{5}b^{10}
$; i d) $x^{8}-8x^{5}+24x^{2}-\frac{32}{x}+\frac{16}{x^{4}}$
\end{solucio}

\begin{exercici}
Sense fer el desenvolupament, troba els termes $2^{o}$ i $14^{o}$ del
desenvolupament de $(x-1)^{14}$.
\end{exercici}
\begin{indicacio}
Usal terme general del binomi de Newton: el terme corresponent depèn del coeficient binomial i de les potències dels dos termes del binomi.
\end{indicacio}

\begin{solucio}
$-14x^{13}$ i $-14x$
\end{solucio}

\begin{exercici}
Sense fer el desenvolupament, troba el terme central del desenvolupament
de $(a^{2}x+x^{2})^{6}$.
\end{exercici}
\begin{indicacio}
Usal terme general del binomi de Newton: el terme corresponent depèn del coeficient binomial i de les potències dels dos termes del binomi.
\end{indicacio}

\begin{solucio}
$20a^{6}x^{9}$
\end{solucio}

\begin{exercici}
Troba el terme independent del desenvolupament de $\left( 3x-\frac{1
}{x}\right) ^{8}$.
\end{exercici}
\begin{indicacio}
Usal terme general del binomi de Newton: el terme corresponent depèn del coeficient binomial i de les potències dels dos termes del binomi.
\end{indicacio}

\begin{solucio}
$5670$
\end{solucio}

\begin{exercici}
Troba el terme la qual part literal és $x^{5}$ en desenvolupar $
\left( x+\frac{2}{x}\right) ^{7}$.
\end{exercici}
\begin{indicacio}
Usal terme general del binomi de Newton: el terme corresponent depèn del coeficient binomial i de les potències dels dos termes del binomi.
\end{indicacio}

\begin{solucio}
$14x^{5}$
\end{solucio}

\section{Divisió de polinomis i algorisme d’Euclides}

\begin{exercici}
Efectua les divisions següents de polinomis:

\begin{enumerate}
\item[a)] $(x^{3}-3x^{2}+2x+5):(x^{2}+x+2)$

\item[b)] $(4x^{5}-3x^{3}+2x+1):(2x^{3}-x+2)$

\item[c)] $\left( \frac{1}{2}x^{6}+3x^{4}-\frac{1}{3}x^{3}+\frac{1}{6}
\right):(3x^{3}+2x)$

\item[d)] $(x^{5}-5x^{4}-3x+15):(x^{2}-\sqrt{3})$

\item[e)] $(x^{3}+2x^{2}-3x-6):(6x+1)$
\end{enumerate}
\end{exercici}
\begin{indicacio}
Aplica la divisió euclidiana de polinomis. En l’algorisme d’Euclides, continua fins que el residu sigui zero.
\end{indicacio}

\begin{solucio}
\[
\begin{array}{lll}
& \text{Quocient} & \text{Residu} \\
a) & x-4 & 4x+13 \\
b) & 2x^{2}-\frac{1}{2} & -4x^{2}+\frac{3}{2}x+2 \\
c) & \frac{1}{6}x^{3}+\frac{8}{9}x-\frac{1}{9} & -\frac{16}{9}x^{2}+\frac{2}{
9}x+\frac{1}{6} \\
d) & x^{3}-5x^{2}+\sqrt{3}x-5\sqrt{3} & 0 \\
e) & \frac{1}{6}x^{2}+\frac{11}{36}x-\frac{119}{216} & -\frac{1177}{216}
\end{array}
\]
\end{solucio}

\begin{exercici}
Troba els valors de $m$ i $n$ perquè el polinomi $p(x)=x^{2}-10$
sigui un divisor de $q(x)=x^{4}-3x^{3}-8x^{2}+mx+n$.
\end{exercici}
\begin{indicacio}
Iguala els coeficients dels termes del mateix grau o aplica el teorema del residu, segons correspongui.
\end{indicacio}

\begin{solucio}
$m=30$ i $n=-20$
\end{solucio}

\begin{exercici}
Troba els valors de $m$ i $n$ perquè el polinomi $
p(x)=x^{4}+4x^{3}-x^{2}+m$ sigui un múltiplo de $q(x)=x^{2}+7x+n$.
\end{exercici}
\begin{indicacio}
Iguala els coeficients dels termes del mateix grau o aplica el teorema del residu, segons correspongui.
\end{indicacio}

\begin{solucio}
$m=84$ i $n=14$
\end{solucio}

\begin{exercici}
Troba el m.c.d. dels polinomis següents per l'algorisme de
Euclides:

\begin{enumerate}
\item[a)] $p(x)=x^{3}+x^{2}+x+1$ i $q(x)=x^{3}-x^{2}+x-1$

\item[b)] $p(x)=x^{4}-5x^{3}+6x^{2}-15x+9$ i $q(x)=x^{3}-7x^{2}+13x-6$

\item[c)] $p(x)=x^{4}-3x^{3}-12x^{2}+17x-3$ i $
q(x)=x^{4}-5x^{3}+3x^{2}-10x+2 $
\end{enumerate}
\end{exercici}
\begin{indicacio}
Aplica la divisió euclidiana de polinomis. En l’algorisme d’Euclides, continua fins que el residu sigui zero.
\end{indicacio}

\begin{solucio}
\[
\begin{array}{ll}
a) & x^{2}+1 \\
b) & x^{2}-5x+3 \\
c) & x^{2}-5x+1
\end{array}
\]
\end{solucio}

\section{Regla de Ruffini, residu i factorització}

\begin{exercici}
Efectua les divisions següents mitjançant la regla de Ruffini:

\begin{enumerate}
\item[a)] $(x^{5}-9x^{3}+30x^{2}):(x+4)$

\item[b)] $(2x^{5}-x^{2}-x-1):(x+\frac{1}{3})$

\item[c)] $(x^{4}+2x^{3}+3x^{2}-4x-10):(x-\sqrt{2})$

\item[d)] $(x^{4}+ax^{3}+a^{2}x^{2}+a^{3}x+a^{4}):(x+a)$
\end{enumerate}
\end{exercici}
\begin{indicacio}
En les divisions per un binomi de primer grau, pots aplicar Ruffini o el teorema del residu.
\end{indicacio}

\begin{solucio}
\[
\begin{array}{lll}
& \text{Quocient} & \text{Residu} \\
a) & x^{4}-4x^{3}+7x^{2}+2x-8 & 32 \\
b) & 2x^{4}-\frac{2}{3}x^{3}+\frac{2}{9}x^{2}-\frac{29}{27}x-\frac{52}{81} &
-\frac{191}{243} \\
c) & x^{3}+(2+\sqrt{2})x^{2}+(5+2\sqrt{2})x+5\sqrt{2} & 0 \\
d) & x^{3}+a^{2}x & a^{4}
\end{array}
\]
\end{solucio}

\begin{exercici}
Troba el residu de les divisions següents: (a) $(x^{128}+1):(x+1)$ i
(b) $(x^{128}-1):(x+1)$.
\end{exercici}
\begin{indicacio}
En les divisions per un binomi de primer grau, pots aplicar Ruffini o el teorema del residu.
\end{indicacio}

\begin{solucio}
(a) $2$\ i (b) $0$
\end{solucio}

\begin{exercici}
Troba el valor de $a$ perquè el polinomi $p(x)=x^{3}-x^{2}+2x-a-3$
sigui divisible per $x+7$.
\end{exercici}
\begin{indicacio}
Iguala els coeficients dels termes del mateix grau o aplica el teorema del residu, segons correspongui.
\end{indicacio}

\begin{solucio}
$a=409$
\end{solucio}

\begin{exercici}
Troba el valor de $k$ perquè en dividir el polinomi $
p(x)=x^{4}-2kx^{3}+x^{2}-4kx+9$ per $x+1$ el residu sigui $-7$.
\end{exercici}
\begin{indicacio}
Iguala els coeficients dels termes del mateix grau o aplica el teorema del residu, segons correspongui.
\end{indicacio}

\begin{solucio}
$k=-3$
\end{solucio}

\begin{exercici}
Troba un polinomi de segon grau que en dividir-lo per $x+1$ el
residu sigui $-2$, que en dividir-lo per $x-2$ el residu sigui $4$, i que sigui mú
ltiplo de $x+3$.
\end{exercici}
\begin{indicacio}
Iguala els coeficients dels termes del mateix grau o aplica el teorema del residu, segons correspongui.
\end{indicacio}

\begin{solucio}
$-\frac{3}{5}x^{2}+\frac{7}{5}x-\frac{6}{5}$
\end{solucio}

\begin{exercici}
Donat el polinomi $p(x)=x^{3}-5x^{2}+ax+b$, troba $a$ i $b$ sabent
que tant en dividir-lo per $x-2$ com per $x+1$ el residu de la división
és $-11$.
\end{exercici}
\begin{indicacio}
Iguala els coeficients dels termes del mateix grau o aplica el teorema del residu, segons correspongui.
\end{indicacio}

\begin{solucio}
$a=2$ i $b=-3$
\end{solucio}

\section{Arrels, màxim comú divisor i mínim comú múltiple}

\begin{exercici}
Troba les arrels i la descomposición factorial dels
següents polinomis:

\begin{enumerate}
\item[a)] $9x^{2}+6x+1$

\item[b)] $25x^{2}-9$

\item[c)] $x^{4}-6x^{2}+9$

\item[d)] $5x^{5}-10\sqrt{2}x^{4}+10x^{3}$

\item[e)] $3x^{4}-10x^{3}+7x^{2}+4x-4$

\item[f)] $x^{4}+2x^{3}-5x^{2}-18x-36$
\end{enumerate}
\end{exercici}
\begin{indicacio}
Factoritza els polinomis tant com sigui possible. Les arrels i els factors lineals són la clau del càlcul.
\end{indicacio}

\begin{solucio}
\[
\begin{array}{lll}
a) & (3x+1)^{2} & -\frac{1}{3}\text{ \ (doble)} \\
b) & (5x+3)(5x-3) & -\frac{3}{5}\text{ \ i }\frac{3}{5} \\
c) & (x-\sqrt{3})^{2}(x+\sqrt{3})^{2} & \pm \sqrt{3}\text{ \ (dobles) } \\
d) & 5x^{3}(x-\sqrt{2})^{2} & 0\text{ \ (triple) i \ }\sqrt{2}\text{ \
(doble)} \\
e) & 3(x-1)^{2}(x-2)(x-2/3) & 1\text{ \ (doble), \ }2\text{ \ i \ }2/3 \\
f) & (x-3)(x+3)(x^{2}+2x+4) & 3\text{ \ i \ }-3
\end{array}
\]
\end{solucio}

\begin{exercici}
Troba el m.c.d. i el m.c.m. dels polinomis següents:

\begin{enumerate}
\item[a)] $4x-4,x^{2}-1$ i $6x^{2}-6$

\item[b)] $4x^{2}-4x+1,2x^{2}-3x+1$ i $x-1$

\item[c)] $x^{3}+5x^{2}-x-5$ i $x^{3}-2x^{2}-x+2$

\item[d)] $x^{6}-1$ i $x^{3}-x$
\end{enumerate}
\end{exercici}
\begin{indicacio}
Factoritza els polinomis tant com sigui possible. Les arrels i els factors lineals són la clau del càlcul.
\end{indicacio}

\begin{solucio}
\[
\begin{array}{lll}
& \text{m.c.d.} & \text{m.c.m.} \\
a) & x-1 & x^{2}-1 \\
b) & 1 & \frac{1}{4}(2x-1)^{2}(x-1) \\
c) & (x-1)(x+1) & \left( x+5\right) \left( x-1\right) \left( x-2\right)
\left( x+1\right) \\
d) & (x-1)(x+1) & x(x^{6}-1)
\end{array}
\]
\end{solucio}

\begin{exercici}
Resol les equacions següents:

\begin{enumerate}
\item[a)] $8x^{4}-38x^{3}+37x^{2}+17x-6=0$

\item[b)] $3x(x+2)-(x+2)(x^{2}-1)=5$
\end{enumerate}
\end{exercici}
\begin{indicacio}
Factoritza els polinomis tant com sigui possible. Les arrels i els factors lineals són la clau del càlcul.
\end{indicacio}

\begin{solucio}
(a) $2,3,1/4$ i $-1/2$; (b) $3,-1+\sqrt{2}$ i $-1-\sqrt{2}$
\end{solucio}

\begin{exercici}
Troba els valors de $m$ i $n$ perquè les fraccions següents siguin
equivalents:
\[
\frac{x^{2}-4}{x^{3}-3x-2}\text{ \ \ \ \ i \ \ \ \ }\frac{x^{2}-x-6}{
x^{3}-x^{2}+mx+n}
\]
\end{exercici}
\begin{indicacio}
Factoritza els polinomis tant com sigui possible. Les arrels i els factors lineals són la clau del càlcul.
\end{indicacio}

\begin{solucio}
$m=-5$ i $n=-3$
\end{solucio}

\section{Fraccions racionals}

\begin{exercici}
Simplifica les fraccions racionals següents:

\begin{enumerate}
\item[a)]
\[
\frac{3x^{2}}{x^{2}-6x}
\]

\item[b)]
\[
\frac{x^{2}-14x+49}{x^{2}-49}
\]

\item[c)]
\[
\frac{x^{4}-16}{x^{4}+8x^{2}+16}
\]

\item[d)]
\[
\frac{x^{3}-3x^{2}+3x-1}{x^{4}+18x^{2}+81}
\]

\item[e)]
\[
\frac{x^{3}-3x+2}{x^{4}+x^{2}+2x-4}
\]
\end{enumerate}
\end{exercici}
\begin{indicacio}
Factoritza numeradors i denominadors, simplifica factors comuns i recorda les restriccions del domini.
\end{indicacio}

\begin{solucio}
\[
\begin{array}{ll}
a) & \frac{3x}{x-6} \\
b) & \frac{x-7}{x+7} \\
c) & \frac{x^{2}-4}{x^{2}+4} \\
d) & \text{Irreductible} \\
e) & \frac{x-1}{x^{2}+2}
\end{array}
\]
\end{solucio}

\begin{exercici}
Efectua les operacions indicades, simplificant el resultat:

\begin{enumerate}
\item[a)]
\[
\frac{1}{x^{2}-4}-\frac{1}{x^{2}-2x}+\frac{1}{x^{2}-4x+4}
\]

\item[b)]
\[
\left( 1-\frac{9}{x^{2}}\right) \cdot \frac{4x}{x^{2}-6x+9}
\]

\item[c)]
\[
\frac{x^{2}-3x}{x^{2}-4x+4}:\left( \frac{x}{x-2}-\frac{x}{x+2}\right)
\]

\item[d)]
\[
\left( \frac{2x}{x^{2}-4}-\frac{x}{x-2}\right):\left( \frac{5}{2+x}-\frac{5
}{2-x}\right)
\]

\item[e)]
\[
\left( \frac{3x}{x-1}+\frac{4}{x+1}-\frac{2x^{2}+7x}{x^{2}-1}\right):\left(
\frac{2x-1}{x-1}-\frac{9}{x+1}\right)
\]

\item[f)]
\[
\frac{x+6}{x^{3}-9x}+\frac{18}{x^{4}-81}+\frac{x+6}{x^{3}+9x}
\]
\end{enumerate}
\end{exercici}
\begin{indicacio}
Factoritza numeradors i denominadors, simplifica factors comuns i recorda les restriccions del domini.
\end{indicacio}

\begin{solucio}
\[
\begin{array}{ll}
a) & \frac{x^{2}+4}{x(x+2)(x-2)^{2}} \\
b) & \frac{4x+12}{x^{2}-3x} \\
c) & \frac{-x^{2}+x+6}{4x-8} \\
d) & \frac{4-x}{10} \\
e) & \frac{x+2}{2x-4} \\
f) & \frac{2x+6}{x^{3}-3x^{2}+9x-27}
\end{array}
\]
\end{solucio}

\begin{exercici}
Troba $A,B$ i $C$, de manera que es compleixin les igualtats que
s'indiquen:

\begin{enumerate}
\item[a)]
\[
\frac{4x-12}{x^{2}-36}=\frac{A}{x-6}+\frac{B}{x+6}
\]

\item[b)]
\[
\frac{x^{2}-2x-1}{x^{2}(x-1)}=\frac{A}{x}+\frac{B}{x^{2}}+\frac{C}{x-1}
\]

\item[c)]
\[
\frac{3x+2}{x^{3}+4x}=\frac{A}{x}+\frac{Bx+C}{x^{2}+4}
\]

\item[d)]
\[
\frac{2x-1}{x^{2}-6x+9}=\frac{A}{x-3}+\frac{B}{(x-3)^{2}}
\]
\end{enumerate}
\end{exercici}
\begin{indicacio}
Factoritza numeradors i denominadors, simplifica factors comuns i recorda les restriccions del domini.
\end{indicacio}

\begin{solucio}
\[
\begin{array}{llll}
& A & B & C \\
a) & 1 & 3 & \\
b) & 3 & 1 & -2 \\
c) & 1/2 & -1/2 & 3 \\
d) & 2 & 5 &
\end{array}
\]
\end{solucio}

\begin{exercici}
Representa la funció $f(x)=x^{2}-6x-7$ i respon les preguntes
següents: (a) Quin és el recorregut d'aquesta funció?
(b) Quines són les antiimatges de $0$ i $-20$?
(c) Per a quins valors de $x$ es compleix $f(x)<0$?
I $f(x)\leq -12$?
\end{exercici}
\begin{indicacio}
Factoritza numeradors i denominadors, simplifica factors comuns i recorda les restriccions del domini.
\end{indicacio}

\begin{solucio}
(a) $[-16,+\infty )$. (b) Les antiimatges de $0$ són $7$ i $-1$; $-20$ no té antiimatges. (c) $f(x)<0$ si $x\in (-1,7)$ i $f(x)\leq
-12$ si $x\in [ 1,5]$.
\begin{center}
\includegraphics[width=0.55\textwidth]{poli31.png}
\end{center}
\end{solucio}

\section{Funcions polinòmiques i racionals}

\begin{exercici}
Representa la funció $f(x)=3x-2$ i respon les preguntes
següents: (a) Quines són les imatges de $1/3$ i
$0$? (b) Quines són les antiimatges de $-1$ i $2$?
(c) Per a quins valors de $x$ es compleix $f(x)\geq 0$?
I $f(x)<4$?
\end{exercici}
\begin{indicacio}
Identifica els punts de tall amb els eixos, el vèrtex o els zeros de la funció, segons el tipus de funció.
\end{indicacio}

\begin{solucio}
(a) $f(1/3)=-1$ i $f(0)=-2$; (b) $1/3$ és antiimatge de $-1$ i $4/3$ ho és de $2$; (c) $f(x)\geq 0$ si $x\in [ 2/3,+\infty )$ i $
f(x)<4$ si $x\in (-\infty,2)$.
\begin{center}
\includegraphics[width=0.55\textwidth]{poli32.png}
\end{center}
\end{solucio}

\begin{exercici}
La gràfica de la funció $f(x)=3x^{2}+bx+4$ passa pel punt $
(2,28)$. Troba el valor de $b$ i representa la gràfica de la funció.
\end{exercici}
\begin{indicacio}
Identifica els punts de tall amb els eixos, el vèrtex o els zeros de la funció, segons el tipus de funció.
\end{indicacio}

\begin{solucio}
$b=6$
\begin{center}
\includegraphics[width=0.55\textwidth]{poli36.png}
\end{center}
\end{solucio}

\begin{exercici}
Escriu una funció polinòmica de tercer grau que tingui per
zeros els nombres $-3,0$ i $4$. Representa aproximadament aquesta funció.
\end{exercici}
\begin{indicacio}
Identifica els punts de tall amb els eixos, el vèrtex o els zeros de la funció, segons el tipus de funció.
\end{indicacio}

\begin{solucio}
$f(x)=x(x+3)(x-4)$
\begin{center}
\includegraphics[width=0.55\textwidth]{poli33.png}
\end{center}
\end{solucio}

\begin{exercici}
Escriu una funció polinòmica que tingui per gràfica la de
la figura següent.\begin{center}
\includegraphics[width=0.48\textwidth]{poli34.png}
\end{center}
\end{exercici}
\begin{indicacio}
Identifica els punts de tall amb els eixos, el vèrtex o els zeros de la funció, segons el tipus de funció.
\end{indicacio}

\begin{solucio}
La funció té 4 zeros: $-2,1,3$ i $5$. Aleshores, $f(x)=k(x+2)(x-1)(x-3)(x-5)$. A més, com que passa per $(0,-2)$, aleshores $k=1/15$. Per tant, la funció és $f(x)=1/15\cdot (x+2)(x-1)(x-3)(x-5)$
\end{solucio}

\begin{exercici}
És possible que la gràfica d'una funció
polinòmica de tercer grau tingui exactament dos punts de tall amb el
eix d'abscisses?
\end{exercici}
\begin{indicacio}
Factoritza el numerador i estudia el signe de cada factor i del denominador en els intervals determinats pels punts crítics.
\end{indicacio}

\begin{solucio}
Sí, per exemple, la funció $f(x)=(x+2)(x-1)^{2}$
té un punt doble en $x=1$. Geomètricament, la funció és tangent a l'eix d'abscisses en $x=1$.
\begin{center}
\includegraphics[width=0.55\textwidth]{poli35.png}
\end{center}
\end{solucio}

\begin{exercici}
Troba els intervals en què la següent funció racional
\[
f(x)=\frac{3x^{2}-5x+2}{2x+9}
\]
és positiva.
\end{exercici}
\begin{indicacio}
Resol pas a pas i comproval resultat substituint-lo en l’enunciat.
\end{indicacio}

\begin{solucio}
Es té
\[
\frac{3x^{2}-5x+2}{2x+9}>0
\]
quan $x\in \left( -9/2,2/3\right) \cup (1,+\infty )$
\end{solucio}
